\bloque{derivadas-integrales}
\tema{variacion}{Derivada y variación}
\dificultad{1}
\section{Interpretación de la derivada}
La derivada describe cómo cambia una función cerca de un punto de su dominio.
Su interpretación geométrica es la pendiente de la recta tangente cuando ese
concepto está definido. Para aproximarla se comparan imágenes de entradas próximas
y se divide su diferencia entre la diferencia de las entradas correspondientes.
Al acercar las entradas se estudia el límite de esas pendientes de rectas secantes.
Un valor positivo de la derivada indica crecimiento local y un valor negativo
indica decrecimiento local, bajo las condiciones habituales de diferenciabilidad.
Una derivada nula por sí sola no permite concluir que exista un máximo o mínimo.
Por ejemplo, para $f(x)=x^{2}$ se obtiene $f'(x)=2\cdot x$.
En este ejemplo la pendiente cambia de signo alrededor del origen.
Conviene distinguir el valor de la función y el valor de su derivada:
uno describe una altura y el otro una tasa de cambio. Sus unidades también
pueden ser distintas. Para interpretar una gráfica, identifica los ejes, selecciona
un punto y describe con palabras qué representa cada valor antes de operar.
